% !TeX root = ../概率论与数理统计笔记.tex
% 由主文档通过 \input 载入；请编译主文档。
\setcounter{chapter}{2} % 保留原书章号。
\chapter{随机变量、分布函数与常见分布}
\label{chap:ps-univariate}

本章把事件模型转化为数值模型，建立描述单个随机量的统一方法。
先修内容为事件与条件概率；密度部分需要一元积分。学习目标是从试验机制识别分布，并按分布计算事件概率。

\section{随机变量与分布函数}
\label{sec:ps-random-variable-cdf}
定义实值\emph{随机变量}（random variable）$X$，说明对每个实数 $x$，集合 $\{X\leq x\}$ 都必须是事件。
定义\emph{分布函数}（cumulative distribution function, CDF）$F_X(x)=\mathbb P(X\leq x)$，证明其单调性、右连续性和两端极限性质。
由分布函数的跳跃计算单点概率，并用左极限处理开闭端点。

\section{离散型随机变量与常见计数模型}
\label{sec:ps-discrete-distributions}
定义\emph{概率质量函数}（probability mass function, PMF），说明非负性与总和为一。
依次介绍伯努利分布（Bernoulli）、二项分布（binomial）、超几何分布（hypergeometric）与泊松分布（Poisson），列出参数范围、取值范围和试验机制。
比较固定次数独立试验、无放回抽样与稀少事件计数；泊松近似的极限条件在第\ref{chap:ps-limits}章讨论。

\section{等待型离散分布与记号约定}
\label{sec:ps-discrete-waiting}
介绍几何分布（geometric）与负二项分布（negative binomial），统一以“获得指定成功次数所需的总试验次数”为基本版本。
另列“成功前的失败次数”版本及其平移关系，防止取值起点和参数化混用。
由独立伯努利试验推导质量函数，并证明几何分布的无记忆性。

\section{密度与常见连续分布}
\label{sec:ps-continuous-distributions}
定义\emph{概率密度函数}（probability density function, PDF），明确本书的连续型计算主要指具有密度的绝对连续分布。
依次介绍均匀分布（uniform）、指数分布（exponential）、正态分布（normal）、伽马分布（gamma）与贝塔分布（beta）。
伽马分布统一采用形状与率参数化，同时给出尺度与率的倒数关系；密度值不直接等于单点概率。

\section{分位数、混合分布与分布的边界}
\label{sec:ps-quantiles-mixtures}
对 $0<u<1$ 定义\emph{分位数}（quantile）$q_u=\inf\{x:F_X(x)\geq u\}$，说明非严格单调分布函数仍可使用广义逆。
用“零点有质量、正半轴有密度”的模型介绍\emph{混合分布}（mixed distribution）。
指出分布函数连续并不必然意味着存在密度，离散与具有密度两类也不穷尽一切分布。

\section{例题、习题与学习检查}
规划按分段分布函数求概率、按抽样机制选择二项或超几何模型、按等待机制选择几何或指数模型、正态标准化与分位数计算。
习题要求先检查质量函数或密度是否归一化，再计算；反例题区分“同一分布”与“同一个随机变量”。
本章暂不集中计算均值和方差，相关数字特征在第\ref{chap:ps-moments}章统一建立。






